\subsectionaddtolist{الف}{\lr{Domain Coupling}}
این مفهوم به میزان وابستگی و ارتباط متقابل بین بخش‌ها یا دامنه‌های مختلف یک سیستم نرم‌افزاری اشاره دارد. به زبان ساده، وقتی دو بخش از سیستم شدیدا به هم متصل باشند (\lr{Coupling} بالا)، تغییر در یکی باعث خرابی یا نیاز به تغییر در دیگری می‌شود. در معماری‌های نوین (مثل میکروسرویس‌ها)، هدف این است که این وابستگی را به حداقل برسانند تا هر بخش بتواند به صورت مستقل توسعه یافته و تغییر کند.

\subsectionaddtolist{ب}{\lr{Data Decomposition} }
تجزیه داده‌ها به معنای شکستن و تقسیم کردن یک پایگاه داده بزرگ و یکپارچه به دیتابیس‌های کوچک‌تر و مجزا است؛ به طوری که هر میکروسرویس یا دامنه فقط به داده‌های مربوط به خودش دسترسی داشته باشد. به زبان ساده، یعنی به جای اینکه همه بخش‌های یک سیستم اطلاعات خود را در یک انبار مشترک بریزند، هر بخش انبار اختصاصی کوچک خودش را داشته باشد تا تداخلی در کار هم ایجاد نکنند.

\subsectionaddtolist{ج}{\lr{Event-Driven Communication Pattern}}
یک روش ارتباطی بین سرویس‌های مختلف است که در آن، اجزای سیستم از طریق ارسال و دریافت رویدادها با یکدیگر صحبت می‌کنند. به زبان ساده، وقتی در یک بخش از سیستم اتفاقی می‌افتد (مثلا خرید کاربر ثبت شد)، آن بخش یک اعلان یا سیگنال به هوا می‌فرستد. سرویس‌های دیگر که به این موضوع علاقه دارند، این سیگنال را دریافت کرده و کار خودشان را انجام می‌دهند، بدون اینکه نیاز باشد مستقیما با سرویس اولیه تماس بگیرند یا منتظر پاسخ آن بمانند.
